% 分类目录：ml
\documentclass[12pt, a4paper]{article}
\usepackage[margin=2.5cm]{geometry}
\usepackage{ctex}
\usepackage{amsmath, amssymb}
\usepackage{listings}
\usepackage{xcolor}
\usepackage{tabularx}
\usepackage{hyperref}

\lstset{
    language=Python,
    basicstyle=\ttfamily\small,
    keywordstyle=\color{blue},
    commentstyle=\color{green!50!black},
    stringstyle=\color{red},
    showstringspaces=false,
    breaklines=true,
    numbers=left,
    numberstyle=\tiny\color{gray},
    frame=single
}

\title{知识点：聚类方法 (Clustering Methods)}
\author{}
\date{}

\begin{document}
\maketitle

\section{核心定义}
\textbf{中文定义}：聚类（Clustering）是机器学习中最重要的无监督学习任务之一。其目标是将数据集划分为若干个不相交的子集（称为“簇”或“类”），使得同一个簇内的样本尽可能相似，而不同簇之间的样本尽可能不同。

\textbf{英文定义}：Clustering is one of the most fundamental unsupervised learning tasks. Its objective is to partition a dataset into several disjoint subsets (called "clusters"), such that samples within the same cluster are as similar as possible, while samples from different clusters are as distinct as possible.

\section{划分式聚类 (Partitioning Methods)}

\subsection{K-Means (K-均值算法)}
\textbf{原理}：通过迭代寻找 $K$ 个簇的中心，并根据样本到中心的距离进行切分。
\begin{enumerate}
    \item 随机初始化 $K$ 个质心。
    \item 将每个点分配给距离最近的质心。
    \item 重新计算每个簇的质心（均值）。
    \item 重复直到收敛（质心不再变化或达到最大迭代次数）。
\end{enumerate}
\textbf{优化}：$K$ 值的选择常用“手肘法”（Elbow Method）或轮廓系数。

\subsection{K-Medoids (K-中心点)}
相比于均值，K-Medoids 选择簇中真实的样本作为中心，对异常值（Outliers）具有更好的鲁棒性。

\section{层次聚类 (Hierarchical Methods)}

\subsection{AGNES (Agglomerative Nesting)}
一种“自底向上”的聚类策略。
\begin{itemize}
    \item 最初将每个样本视为一个簇。
    \item 每次合并距离最近的两个簇。
    \item 距离度量（Linkage）：最小距离（Single-linkage）、最大距离（Complete-linkage）、平均距离（Average-linkage）。
\end{itemize}

\section{密度聚类 (Density-based Methods)}

\subsection{DBSCAN}
\textbf{基于密度的空间聚类算法}。
\begin{itemize}
    \item \textbf{核心参数}：$\epsilon$（邻域半径）和 $MinPts$（最小样本数）。
    \item \textbf{概念}：核心点、边界点、噪声点（离群点）。
    \item \textbf{优点}：可以发现任意形状的簇，且能自动识别噪声点，无需预先指定簇数。
\end{itemize}

\section{模型聚类 (Model-based Methods)}

\subsection{高斯混合模型 (GMM)}
\textbf{原理}：假设数据是由多个高斯分布组合而成的模型产生的。
\begin{itemize}
    \item 属于一种“软聚类”（Soft Clustering），即每个样本以一定的概率属于某个簇。
    \item \textbf{求解}：使用期望最大化算法（EM Algorithm）进行迭代求解。
\end{itemize}

\section{聚类性能评估}

\subsection{内部指标 (Internal Metrics)}
无需参考真实标签。
\begin{itemize}
    \item \textbf{轮廓系数 (Silhouette Coefficient)}：取值范围 $[-1, 1]$，接近 1 表示聚类效果好。
    \item \textbf{DBI (Davies-Bouldin Index)}：值越小表示簇内越紧广、簇间越分散。
\end{itemize}

\subsection{外部指标 (External Metrics)}
需要参考真实标签（如评估算法性能）。
\begin{itemize}
    \item \textbf{ARI (Adjusted Rand Index)}：衡量聚类结果与真实标签的匹配程度。
    \item \textbf{NMI (Normalized Mutual Information)}：归一化互信息。
\end{itemize}

\section{算法对比总结}
\begin{table}[h]
\centering
\begin{tabularx}{\textwidth}{|l|X|X|}
\hline
\textbf{算法} & \textbf{核心优势} & \textbf{典型缺点} \\
\hline
K-Means & 速度快、可扩展性强 & 需要预设 $K$、对噪声敏感、只能发现凸形簇 \\
\hline
Hierarchical & 可解释性强（树状图） & 计算复杂度高 ($O(n^3)$) \\
\hline
DBSCAN & 识别任意形状、抗噪 & 对参数 $\epsilon$ 极其敏感、难以处理密度不均的数据 \\
\hline
GMM & 概率输出（软分配） & 计算量大、容易陷入局部最优 \\
\hline
\end{tabularx}
\end{table}

\section{关键代码片段 (Python)}
\begin{lstlisting}
from sklearn.cluster import KMeans, DBSCAN
from sklearn.mixture import GaussianMixture
from sklearn.metrics import silhouette_score

# 1. K-Means
kmeans = KMeans(n_clusters=3, random_state=42)
labels_k = kmeans.fit_predict(X)

# 2. DBSCAN
dbscan = DBSCAN(eps=0.5, min_samples=5)
labels_db = dbscan.fit_predict(X)

# 3. 评估
score = silhouette_score(X, labels_k)
print(f"Silhouette Score: {score}")
\end{lstlisting}

\section{深度面试题}

\textbf{问：K-Means 初始质心选择不当会产生什么后果？如何改进？}\\
答：可能导致算法陷入局部最优解，或增加迭代次数。改进方法是使用 \textbf{K-Means++} 算法，其核心思想是让初始质心之间的距离尽可能远，从而加速收敛并获得更好的全局结果。

\textbf{问：聚类算法中的“距离”一定要用欧氏距离吗？}\\
答：不一定。距离的选择取决于数据特征和业务场景。例如，文档聚类常用余弦相似度，离散型特征可用 Jaccard 距离，带有权重的场景可用马氏距离。

\section{参考文献}
\begin{enumerate}
    \item Lloyd, S. P. (1982). \textit{Least Squares Quantization in PCM}. IEEE Transactions on Information Theory.
    \item Ester, M., et al. (1996). \textit{A Density-Based Algorithm for Discovering Clusters in Large Spatial Databases with Noise}. KDD.
    \item Reynolds, D. A. (2009). \textit{Gaussian Mixture Models}. Encyclopedia of Biometric.
\end{enumerate}

\end{document}
